\documentclass[11pt,a4paper]{article}
\usepackage[utf8]{inputenc}
\usepackage[T1]{fontenc}
\usepackage[english]{babel}
\usepackage{geometry}
\usepackage{xcolor}
\usepackage{graphicx}
\usepackage{tikz}
\usepackage{fancyhdr}
\usepackage{titlesec}
\usepackage{hyperref}
\usepackage{setspace}
\usepackage{enumitem}
\usepackage{array}
\usepackage{booktabs}
\usepackage{colortbl}
\usepackage{amsmath}
\usepackage{subcaption}
\usepackage{float}

% ── Colours ───────────────────────────────────────────────────────────────────
\definecolor{bleuclair}{HTML}{E3F2FD}
\definecolor{chaptercolor}{HTML}{1A5490}
\definecolor{sectioncolor}{HTML}{2563EB}
\definecolor{subsectioncolor}{HTML}{7C3AED}
\definecolor{tableheader}{HTML}{1A5490}
\definecolor{tablerowA}{HTML}{EBF3FB}

% ── Section formatting ────────────────────────────────────────────────────────
\titleformat{\section}{\normalfont\large\bfseries\color{sectioncolor}}{}{0pt}{}[\vspace{-4pt}\color{sectioncolor}\hrule\vspace{3pt}]
\titleformat{\subsection}{\normalfont\normalsize\bfseries\color{subsectioncolor}}{}{0pt}{}
\titlespacing{\section}{0pt}{8pt}{4pt}
\titlespacing{\subsection}{0pt}{5pt}{2pt}

% ── Page layout ───────────────────────────────────────────────────────────────
\geometry{top=2cm, bottom=2cm, left=2.5cm, right=2.5cm}
\setlength{\headheight}{14pt}
\pagestyle{fancy}
\fancyhf{}
\fancyhead[L]{\small Computer Vision}
\fancyhead[R]{\small Project 4 --- Panorama Stitcher}
\fancyfoot[C]{\thepage}
\renewcommand{\headrulewidth}{0.5pt}
\setstretch{1.15}
\hypersetup{colorlinks=true, linkcolor=sectioncolor, urlcolor=blue}
\graphicspath{{./}}

\begin{document}

%% ═══════════════════════════════════════════════════════════════ COVER PAGE ══
% ==================== PAGE DE GARDE ====================
\begin{titlepage}

    \begin{tikzpicture}[remember picture, overlay]

        \fill[bleuclair] (current page.south west) rectangle (current page.north east);

        \node[opacity=0.1] at (current page.center) {\includegraphics[width=20cm]{logo-uir.png}};

    \end{tikzpicture}

    \centering

    

    \includegraphics[height=4.5cm]{logo-uir.png}

    \vspace{1cm}

    

    {\color{chaptercolor}\rule{\textwidth}{2pt}}\\[0.6cm]

    {\Huge\bfseries\color{chaptercolor} Rapport de Projet}\\[8pt]

    {\huge\bfseries Oracle Panorama Image Stitching}\\[6pt]


    {\color{chaptercolor}\rule{\textwidth}{2pt}}\\[1.5cm]
    
    \begin{minipage}{0.45\textwidth}
        \centering
        {\Large\textbf{Aouami Salma}}\\[4pt]
        {\large Student ID: 127443}
    \end{minipage}
    \hfill
    \begin{minipage}{0.45\textwidth}
        \centering
        {\Large\textbf{Zineb Arrami}}\\[4pt]
        {\large Student ID: 118789}
    \end{minipage}
    
    \vspace{0.8cm}
    
    \begin{minipage}{0.45\textwidth}
        \centering
        {\Large\textbf{Hiba Arbaoui}}\\[4pt]
        {\large Student ID: 120401}
    \end{minipage}
    \hfill
    \begin{minipage}{0.45\textwidth}
        \centering
        {\Large\textbf{Ayyadi Marwa}}\\[4pt]
        {\large Student ID: 120278}
    \end{minipage}
    
    \vspace{1cm}
    
    \begin{minipage}{0.45\textwidth}
        \begin{flushleft}
            \textbf{Encadré par :}\\
            Ilias Tougui
        \end{flushleft}
    \end{minipage}
    \hfill
    \begin{minipage}{0.45\textwidth}
        \begin{flushright}
            \textbf{Module :}\\
            Computer Vision
        \end{flushright}
    \end{minipage}
    
    \vfill
    {\large\textbf{Année Universitaire : 2025 -- 2026}}
\end{titlepage}


\newpage

%% ═══════════════════════════════════════════════════════════════════ SECTION 1
\section{Overview}

This project builds a fully classical image stitching pipeline combining overlapping photographs into seamless panoramas using SIFT, RANSAC homography estimation, projective warping, and multi-scale blending — no learning-based component is used. Three real datasets were tested: a 2-image coastal sunset (\texttt{data3}), a 3-image elevated panorama (\texttt{data1}), and a 3-image Lac d'Annecy lakeside scene (\texttt{data2}).

%% ═══════════════════════════════════════════════════════════════════ SECTION 2
\section{Pipeline}

\subsection{Preprocessing}
Images are resized to a 1\,500\,px maximum dimension. CLAHE is applied in the L channel of LAB colour space to locally boost contrast in flat regions (sky, water), significantly increasing keypoint yield in otherwise featureless zones.

\subsection{Keypoint Detection --- SIFT}
SIFT builds a Difference-of-Gaussians scale-space pyramid and detects extrema stable across scale and rotation. Each keypoint is described by a \textbf{128-d} L2-normalised gradient histogram vector over a $16{\times}16$ neighbourhood. We use 5\,000--8\,000 keypoints with a lowered contrast threshold (0.01--0.02) to handle low-texture regions.

\subsection{Feature Matching --- FLANN + Lowe Ratio Test}
FLANN with randomised kd-trees retrieves the 2 nearest neighbours per descriptor. Lowe's ratio test accepts a match only if $d_1 < 0.75\,d_2$, removing ambiguous correspondences before homography estimation.

\subsection{Homography Estimation --- RANSAC}
A $3{\times}3$ projective homography $\mathbf{H}$ (8 DoF) is estimated via DLT. RANSAC samples 4 correspondences per iteration (10\,000 iterations, $\varepsilon = 5$\,px inlier threshold). Pairwise transforms are chained to a global canvas frame:
\[
H_i^{\text{canvas}} = H_{i-1}^{\text{canvas}} \cdot H_{i \to i-1}
\]

\subsection{Warping \& Blending}
Each image is warped via \texttt{cv2.warpPerspective} onto a shared canvas sized by bounding all projected corners. Two blending modes were compared:
\begin{itemize}[noitemsep, topsep=2pt]
    \item \textbf{Alpha blending:} pixel weights from distance transform to image boundary. Fast, but ghosting artefacts appear near seams.
    \item \textbf{Multi-band (5-level Laplacian pyramid):} blends each frequency band with an appropriately smoothed mask, eliminating ghosting and gracefully handling exposure differences.
\end{itemize}

%% ═══════════════════════════════════════════════════════════════════ SECTION 3
\section{Results}

\begin{center}
\small
\begin{tabular}{l r r r r}
\toprule
\rowcolor{tableheader}
\color{white}\textbf{Pair} & \color{white}\textbf{Matches} & \color{white}\textbf{Inliers} & \color{white}\textbf{Inlier \%} & \color{white}\textbf{Reproj.\ err.} \\
\midrule
\rowcolor{tablerowA}
data3: 6625 $\leftrightarrow$ 6626 & 505 & 480 & 95.0\% & 0.15\,px \\
data1: 6622 $\leftrightarrow$ 6623 & 281 & 232 & 82.6\% & 1.86\,px \\
\rowcolor{tablerowA}
data1: 6623 $\leftrightarrow$ 6624 & 226 & 201 & 88.9\% & 1.60\,px \\
data2: 6628 $\leftrightarrow$ 6629 & 353 & 266 & 75.4\% & 1.08\,px \\
\rowcolor{tablerowA}
data2: 6629 $\leftrightarrow$ 6630 & 1\,511 & 1\,467 & 97.1\% & 0.89\,px \\
\bottomrule
\end{tabular}
\end{center}

\vspace{0.4em}
\textbf{data3} achieves the lowest error (0.15\,px) due to clean overlap and simple geometry. \textbf{data1} shows higher error (${\approx}1.7$\,px) from slight perspective drift between hand-held portrait shots. \textbf{data2} was the hardest — sparse keypoints on the flat lake required 8\,000 features and threshold 0.01; the textured dock area in pair 6629--6630 compensated with 1\,467 inliers. Multi-band blending eliminated visible ghosting present in all alpha-blended outputs.

\vspace{0.5em}
The three panorama outputs are shown below. All images were blended using the multi-band Laplacian pyramid strategy.

\begin{figure}[H]
    \centering
    \includegraphics[width=0.82\linewidth]{data3_result.png}
    \caption{\textbf{data3} --- Coastal Sunset panorama (2 images, canvas 2\,324\,$\times$\,1\,324\,px). Clean alignment with 95\% inlier rate and sub-pixel reprojection error.}
    \label{fig:data3}
\end{figure}

\begin{figure}[H]
    \centering
    \begin{subfigure}[t]{0.48\linewidth}
        \includegraphics[width=\linewidth]{data2_result.png}
        \caption{\textbf{data2} --- Lac d'Annecy (3 images). Hardest scene due to flat lake surface; CLAHE and higher feature count were essential.}
        \label{fig:data2}
    \end{subfigure}
    \hfill
    \begin{subfigure}[t]{0.48\linewidth}
        \includegraphics[width=\linewidth]{data1_result.png}
        \caption{\textbf{data1} --- Elevated Coastal Panorama (3 images). Portrait-orientation shots introduce slight perspective warp visible at canvas edges.}
        \label{fig:data1}
    \end{subfigure}
    \caption{Multi-band blended panorama outputs for data2 and data1.}
\end{figure}

%% ═══════════════════════════════════════════════════════════════════ SECTION 4
\section{Challenges \& Course Concepts}

The main challenge was \textbf{low-texture regions}: without CLAHE, RANSAC failed entirely on data2. Ratio test tuning was also critical --- below 0.65 too many true matches were rejected; above 0.80 false matches destabilised the homography.

Course concepts applied: \textit{scale-space / DoG} (SIFT), \textit{gradient orientation histograms} (descriptor design), \textit{homogeneous coordinates \& DLT} (homography fitting), \textit{RANSAC} (robust estimation), \textit{projective geometry} (warping), and \textit{Laplacian pyramids} (multi-resolution blending).


\end{document}